\chapter{Introduction}
\label{ch:introduction}

Volumetric data is a collection of samples in 3D space. These samples
represent values of some properties of data obtained using sampling, simulation,
or modeling techniques. A common source of such data is medical imaging with
modalities like \ac{mri}, \ac{ct}, or
\ac{pet}~\cite[Sections~7.1--7.2]{kaufman_7_2005}.

Volume visualization methods are used to extract useful information from
such data. Initially, surface rendering techniques were used, but more
recently volume rendering techniques, most notably path tracing, have been
employed. Because they use the entire 3D structure, they convey more
information but are also more computationally
demanding~\cite[Section~4]{kaufman_43_2000}.

To improve performance, various optimization techniques have been developed.
Some target specific steps of the algorithm, such as empty space
skipping~\cite{hadwiger_sparseleap_2018} or early ray
termination~\cite{matsui_parallel_2005}, while others introduce caching
mechanisms that speed up the global illumination computation
step~\cite{krivanek_radiance_2008}.

As in other fields in recent years, computer graphics and volumetric
rendering have also been shaped by neural networks, from neural importance
sampling~\cite{muller_neural_2019} and caching~\cite{muller_real-time_2021}
for mesh data to deep ambient occlusion~\cite{engel_deep_2021} and
scattering inside clouds~\cite{kallweit_deep_2017} for volumetric data.
Performance is especially critical for interactive volume visualization,
where the image must update in real time in response to user input. This
requires keeping frame times as low as possible, and in many cases neural
approaches have proved more effective than conventional optimization
methods.

Bauer \emph{et~al.}~\cite{bauer_photon_2024} trained a neural network to
predict indirect in-scatter radiance from position, direction, and the phase
function coefficient. Their training data is generated offline using photon
mapping, so the model must be trained in a separate precomputation step that
is repeated for every new volume and delays the first image.

\section{Goal and contributions}

The goal of our work is to shorten the rendering time for generating an image
in volumetric path tracing, both when a volume is first loaded and when the
camera moves, so that the renderer remains responsive and supports interactive
volumetric data analysis.

To achieve this, we propose a new method for online neural caching of indirect
in-scatter radiance that combines and extends existing approaches, runs
on all platforms and in the browser, and is vendor-agnostic. Our thesis
consists of the following contributions:
\begin{itemize}
    \item An online training scheme that generates training data using path
    tracing. Training is done during rendering, removing a separate
    precomputation step and reducing the time to the first image when a new
    volume is loaded.

    \item Learned filtering of radiance values. Path-traced samples are
    generated from the camera view, so we can apply simple 2D image-filtering
    techniques, e.g., bilateral filtering, before passing data to the model.
    The network learns the filtered results, so the filtering cost is incurred
    only during training, not at inference, which is beneficial if the filtering
    step is more expensive.

    \item An open-source, cross-platform, vendor-agnostic implementation
    inside an existing volumetric rendering framework.

    \item An extensive evaluation across different configurations and multiple
    volumes acquired through different techniques. We include all the code for
    our testing framework along with exact experiment definitions to make
    results fully reproducible.
\end{itemize}

\section{Thesis structure}

The remainder of this thesis is organized as follows.
\Cref{ch:related_work} presents notable work on the development of
volumetric rendering, caching approaches, and neural methods, including the latest state of the art.
\Cref{ch:vpt} explains the basic volumetric path tracing preliminaries and
describes the path tracer that serves as the basis for our implementation.
\Cref{ch:methods} details our method, both in terms of the
neural network architecture and the WebGPU implementation.
\Cref{ch:evaluation} presents our evaluation framework, the experiments
we conducted, and the results.
\Cref{ch:conclusion} summarizes our findings and proposes possible directions
for future work.